\textbf{Phần 1. Trắc nghiệm lựa chọn}
\setcounterpageref{ex}{0}

\begin{ex} %[3D4N25]
Phép chia có số dư bằng 0 được gọi là
\choice
{Phép chia có dư}
{Phép nhân}
{\True Phép chia hết}
{Phép chia không thực hiện được}
\loigiai{
Phép chia có số dư bằng 0 là phép chia hết.\\
\textbf{Đáp án của câu hỏi là C}
}
\end{ex}

\begin{ex} %[3D4N25]
Trong phép chia có dư, số dư luôn luôn phải
\choice
{Lớn hơn số chia}
{Bằng số chia}
{Gấp đôi số chia}
{\True Nhỏ hơn số chia}
\loigiai{
Số dư luôn luôn nhỏ hơn số chia.\\
\textbf{Đáp án của câu hỏi là D}
}
\end{ex}

\begin{ex} %[3D4H25]
Kết quả của phép chia $25 : 4$ là
\choice
{\True $6$ dư $1$}
{$5$ dư $5$}
{$6$ dư $2$}
{$7$ dư $1$}
\loigiai{
$25 = 4 \times 6 + 1$ (thương là 6, dư 1).\\
\textbf{Đáp án của câu hỏi là A}
}
\end{ex}

\begin{ex} %[3D4H25]
Kết quả của phép chia $38 : 5$ là
\choice
{$7$ dư $2$}
{\True $7$ dư $3$}
{$8$ dư $3$}
{$6$ dư $8$}
\loigiai{
$38 = 5 \times 7 + 3$ (thương 7, dư 3).\\
\textbf{Đáp án của câu hỏi là B}
}
\end{ex}

\begin{ex} %[3D4H25]
Trong phép chia cho 6, số dư lớn nhất có thể có là
\choice
{$6$}
{$4$}
{$7$}
{\True $5$}
\loigiai{
Số dư lớn nhất luôn kém số chia 1 đơn vị: $6 - 1 = 5$.\\
\textbf{Đáp án của câu hỏi là D}
}
\end{ex}

\begin{ex} %[3D4H25]
Trong phép chia cho 8, số dư lớn nhất có thể có là
\choice
{$8$}
{$6$}
{\True $7$}
{$9$}
\loigiai{
Số dư lớn nhất là 7.\\
\textbf{Đáp án của câu hỏi là C}
}
\end{ex}

\begin{ex} %[3D4V25]
Có 29 cái kẹo chia đều cho 4 bạn. Mỗi bạn được nhiều nhất bao nhiêu cái kẹo và còn thừa mấy cái?
\choice
{$6$ cái, thừa $5$ cái}
{$7$ cái, thừa $2$ cái}
{$8$ cái, thừa $1$ cái}
{\True $7$ cái, thừa $1$ cái}
\loigiai{
$29 : 4 = 7$ dư 1.\\
\textbf{Đáp án của câu hỏi là D}
}
\end{ex}

\begin{ex} %[3D4V25]
Có 45 học sinh đi đò, mỗi chuyến đò chở được nhiều nhất 6 học sinh. Cần ít nhất bao nhiêu chuyến đò để chở hết số học sinh đó?
\choice
{$7$ chuyến}
{\True $8$ chuyến}
{$9$ chuyến}
{$6$ chuyến}
\loigiai{
$45 : 6 = 7$ dư 3 học sinh. Cần thêm 1 chuyến cho 3 bạn còn lại $\Rightarrow 7 + 1 = 8$ chuyến.\\
\textbf{Đáp án của câu hỏi là B}
}
\end{ex}

\begin{ex} %[3D4V25]
Một cuộn dây dài $35\text{ m}$ cắt thành các đoạn dài $4\text{ m}$. Cắt được nhiều nhất số đoạn dây là
\choice
{\True $8$ đoạn (còn thừa $3\text{ m}$)}
{$7$ đoạn}
{$9$ đoạn}
{$6$ đoạn}
\loigiai{
$35 : 4 = 8$ dư 3 $\Rightarrow$ Cắt được 8 đoạn.\\
\textbf{Đáp án của câu hỏi là A}
}
\end{ex}

\begin{ex} %[3D4V25]
Tìm số bị chia trong phép chia có số chia là 7, thương là 8 và số dư là 5.
\choice
{$56$}
{$66$}
{\True $61$}
{$59$}
\loigiai{
Số bị chia $= 7 \times 8 + 5 = 56 + 5 = 61$.\\
\textbf{Đáp án của câu hỏi là C}
}
\end{ex}

\begin{ex} %[3D4V25]
Trong một phép chia có số chia là 6, thương là 12 và số dư là số dư lớn nhất có thể. Số bị chia là
\choice
{$72$}
{\True $77$}
{$78$}
{$75$}
\loigiai{
Số dư lớn nhất là 5. Số bị chia $= 12 \times 6 + 5 = 72 + 5 = 77$.\\
\textbf{Đáp án của câu hỏi là B}
}
\end{ex}

\begin{ex} %[3D4C25]
Có một số cái kẹo, nếu chia đều cho 5 bạn thì thừa 3 cái. Nếu chia đều cho 5 bạn và thêm vào 2 cái kẹo nữa thì mỗi bạn được 8 cái kẹo. Số kẹo ban đầu là
\choice
{\True $38$ cái}
{$36$ cái}
{$40$ cái}
{$43$ cái}
\loigiai{
Tổng số kẹo sau khi thêm: $8 \times 5 = 40$ cái.\\
Số kẹo ban đầu: $40 - 2 = 38$ cái (thỏa mãn $38 : 5 = 7$ dư 3).\\
\textbf{Đáp án của câu hỏi là A}
}
\end{ex}

\textbf{Phần 2. Trắc nghiệm đúng sai}
\setcounterpageref{ex}{0}

\begin{ex} %[3D4H25]
Xét tính đúng sai về phép chia có dư:
\choiceTF
{\True Phép chia $30 : 6 = 5$ là phép chia hết.}
{\True Phép chia $35 : 8 = 4$ dư $3$ là phép chia có dư.}
{Trong phép chia cho 5, số dư có thể bằng 5.}
{\True Trong phép chia cho 5, các số dư có thể là 1, 2, 3, 4.}
\loigiai{
\begin{itemize}
    \item $30 : 6 = 5$: Chia hết. Đúng.
    \item $35 : 8 = 4$ dư 3: Đúng.
    \item Số dư phải nhỏ hơn số chia nên không thể bằng 5. Ý 3 Sai, Ý 4 Đúng.
\end{itemize}
\begin{itemize}
    \item a) \textbf{Đúng}.
    \item b) \textbf{Đúng}.
    \item c) \textbf{Sai}.
    \item d) \textbf{Đúng}.
\end{itemize}
}
\end{ex}

\begin{ex} %[3D4V25]
Có 38 quả cam xếp đều vào các đĩa, mỗi đĩa 5 quả. Xét tính đúng sai:
\choiceTF
{\True Xếp được nhiều nhất 7 đĩa cam.}
{\True Còn thừa lại 3 quả cam.}
{Nếu muốn xếp vừa đủ 8 đĩa đầy thì cần thêm 3 quả cam nữa.}
{\True Nếu muốn xếp vừa đủ 8 đĩa đầy thì cần thêm 2 quả cam nữa.}
\loigiai{
\begin{itemize}
    \item $38 : 5 = 7$ dư 3: Xếp được 7 đĩa. Đúng.
    \item Thừa 3 quả: Đúng.
    \item Để được 8 đĩa cần $8 \times 5 = 40$ quả $\Rightarrow$ cần thêm $40 - 38 = 2$ quả $\neq 3$. Ý 3 Sai, Ý 4 Đúng.
\end{itemize}
\begin{itemize}
    \item a) \textbf{Đúng}.
    \item b) \textbf{Đúng}.
    \item c) \textbf{Sai}.
    \item d) \textbf{Đúng}.
\end{itemize}
}
\end{ex}

\textbf{Phần 3. Trả lời ngắn}
\setcounterpageref{ex}{0}

\begin{ex} %[3D4H25]
Tìm số dư trong phép chia $26 : 4$.
\shortans{2}
\loigiai{
\textbf{Đáp số:} $2$\\
$26 : 4 = 6$ dư 2.
}
\end{ex}

\begin{ex} %[3D4H25]
Trong phép chia cho 9, số dư lớn nhất có thể có là bao nhiêu?
\shortans{8}
\loigiai{
\textbf{Đáp số:} $8$\\
$9 - 1 = 8$.
}
\end{ex}

\begin{ex} %[3D4V25]
Tìm số bị chia trong phép chia có số chia là 6, thương là 9 và số dư là 4.
\shortans{58}
\loigiai{
\textbf{Đáp số:} $58$\\
$6 \times 9 + 4 = 54 + 4 = 58$.
}
\end{ex}

\begin{ex} %[3D4V25]
Có 50 quyển vở chia đều cho 6 bạn. Hỏi còn thừa mấy quyển vở?
\shortans{2}
\loigiai{
\textbf{Đáp số:} $2$\\
$50 : 6 = 8$ dư 2.
}
\end{ex}

\begin{ex} %[3D4V25]
Có 34 bạn học sinh muốn đi taxi, mỗi xe taxi chở được 4 bạn. Hỏi cần ít nhất bao nhiêu xe taxi để chở hết số bạn đó?
\shortans{9}
\loigiai{
\textbf{Đáp số:} $9$\\
$34 : 4 = 8$ dư 2 $\Rightarrow$ Cần ít nhất $8 + 1 = 9$ xe.
}
\end{ex}

\begin{ex} %[3D4C25]
Một số chia cho 8 được thương là 7 và có số dư là số dư lớn nhất. Tìm số đó.
\shortans{63}
\loigiai{
\textbf{Đáp số:} $63$\\
Số dư lớn nhất là 7. Số đó là: $8 \times 7 + 7 = 56 + 7 = 63$.
}
\end{ex}
